\documentclass[11pt,letterpaper]{article}

% ============================================================
% PORTFOLIO DIAGNOSTICS — Research Report #001
% Informe institucional (no paper adaptado)
% ============================================================

\usepackage[T1]{fontenc}
\usepackage[utf8]{inputenc}
\usepackage{lmodern}
\usepackage[spanish,es-tabla,es-nodecimaldot]{babel}
\usepackage[
  top=0.85in,
  bottom=0.85in,
  left=0.95in,
  right=0.95in
]{geometry}
\usepackage{graphicx}
\usepackage{booktabs}
\usepackage{amsmath,amssymb}
\usepackage{array}
\usepackage{caption}
\usepackage{fancyhdr}
\usepackage{titlesec}
\usepackage{microtype}
\usepackage{xcolor}
\usepackage{enumitem}
\usepackage{setspace}
\usepackage{float}
\usepackage[hidelinks]{hyperref}

\graphicspath{{../figures/}}

\definecolor{Navy}{RGB}{27,79,114}
\definecolor{Slate}{RGB}{93,109,126}
\definecolor{Ink}{RGB}{28,40,51}
\definecolor{Muted}{RGB}{110,122,135}
\definecolor{Line}{RGB}{213,216,220}
\definecolor{Soft}{RGB}{245,248,250}
\definecolor{Accent}{RGB}{146,43,33}

\captionsetup{
  font={small,color=Slate},
  labelfont={bf,color=Navy},
  labelsep=period,
  skip=6pt,
  justification=raggedright,
  singlelinecheck=false
}

\titleformat{\section}
  {\color{Navy}\bfseries\Large\raggedright}
  {}{0em}{}
\titlespacing*{\section}{0pt}{1.4em}{0.45em}

\setlist[itemize]{leftmargin=1.2em, itemsep=0.35em, topsep=0.4em, parsep=0pt}

\setlength{\parskip}{0.55em}
\setlength{\parindent}{0pt}
\setstretch{1.2}

\pagestyle{fancy}
\fancyhf{}
\renewcommand{\headrulewidth}{0.4pt}
\renewcommand{\headrule}{\color{Line}\hrule height 0.4pt width\headwidth}
\fancyhead[L]{\small\color{Muted} Villar Capital}
\fancyhead[R]{\small\color{Muted} Informe de investigación nro.\ 001}
\fancyfoot[L]{\small\color{Muted} Diagnóstico de portafolios $|$ Julio 2026}
\fancyfoot[R]{\small\color{Muted}\thepage}
\fancypagestyle{plain}{%
  \fancyhf{}%
  \renewcommand{\headrulewidth}{0pt}%
  \fancyfoot[R]{\small\color{Muted}\thepage}%
}

\newcommand{\kicker}[1]{%
  {\sffamily\small\bfseries\color{Navy}\MakeUppercase{#1}}\par\vspace{0.35em}%
}
\newcommand{\leadq}[1]{%
  {\Large\bfseries\color{Ink}#1}\par\vspace{0.65em}%
}
\newcommand{\insight}[1]{%
  \vspace{0.35em}%
  \noindent\colorbox{Soft}{%
    \parbox{\dimexpr\textwidth-2\fboxsep-1pt\relax}{%
      \vspace{0.45em}%
      {\sffamily\bfseries\color{Navy} Idea clave}\par\vspace{0.2em}%
      {\large\color{Ink}#1}%
      \vspace{0.4em}%
    }%
  }\par\vspace{0.4em}%
}
\newcommand{\fighero}[2]{%
  \begin{figure}[H]
    \centering
    \includegraphics[width=\textwidth,height=0.62\textheight,keepaspectratio]{#1}
    \caption{#2}
  \end{figure}
}
\newcommand{\figfull}[2]{%
  \begin{figure}[H]
    \centering
    \includegraphics[width=0.95\textwidth,height=0.48\textheight,keepaspectratio]{#1}
    \caption{#2}
  \end{figure}
}
\newcommand{\figbig}[2]{%
  \begin{figure}[H]
    \centering
    \includegraphics[width=0.88\textwidth,height=0.44\textheight,keepaspectratio]{#1}
    \caption{#2}
  \end{figure}
}

\begin{document}

% ============================================================
% COVER — mínima: Villar Capital + título + número de reporte
% ============================================================
\begin{titlepage}
\thispagestyle{empty}

\vfill

\begin{center}
{\sffamily\Large\bfseries\color{Navy} Villar Capital}

\vspace{1.1em}
{\sffamily\small\color{Muted} Informe de investigación nro.\ 001}

\vspace{2.8em}
{\fontsize{30}{36}\selectfont\bfseries\color{Ink}%
Más allá del número de activos\par}
\end{center}

\vfill

\noindent{\color{Navy}\rule{\textwidth}{2.5pt}}
\end{titlepage}

% ============================================================
% EXECUTIVE SUMMARY
% ============================================================
\newpage
\kicker{Resumen ejecutivo}
\section*{Lo que revela el diagnóstico}

Una cartera clásica 60/40 con diez instrumentos parece ampliamente diversificada.
Al abrirla capa por capa, la estructura de riesgo cuenta otra historia.

\vspace{1.0em}
\begin{center}
{\LARGE\color{Ink}
$10 \;\rightarrow\; 9{,}55 \;\rightarrow\; 6{,}58 \;\rightarrow\; 3$
\par}
\vspace{0.45em}
{\sffamily\small\color{Muted}
activos \quad capital \quad riesgo \quad fuentes dominantes
}
\end{center}

\vspace{1.0em}
\insight{Contar posiciones sobreestima la diversificación. El capital está disperso; el riesgo, no.}

\vspace{0.8em}
\noindent{\sffamily\small\bfseries\color{Navy} Parámetros del diagnóstico}
\par\vspace{0.45em}
{\sffamily\small\color{Ink}%
\begin{tabular}{@{}p{0.42\textwidth}p{0.52\textwidth}@{}}
Fecha utilizada para la prueba & 2016-01-01 a 2026-06-30 \\[0.25em]
Estructura de la cartera & Pesos objetivo constantes (60/40) \\[0.25em]
Serie de retornos & Logarítmicos diarios \\[0.25em]
Herramientas de dependencia & Correlación y red de agrupamientos \\[0.25em]
Descomposición del riesgo & Contribución al riesgo (Euler) \\[0.25em]
Concentración & Posiciones efectivas (capital y riesgo) \\[0.25em]
Modelo de factores comunes & Análisis de componentes principales \\
\end{tabular}}

% ============================================================
% THE HOOK — Una cartera que cualquiera compraría
% ============================================================
\newpage
\kicker{El punto de partida}
\section*{Una cartera que cualquiera compraría}

Imaginemos un inversor que construye una cartera clásica 60/40.

Diez instrumentos. Siete de renta variable ---índices, tecnología, acciones individuales---.
Tres de renta fija. Dentro de cada bloque, pesos similares.

A simple vista, parece una cartera ampliamente diversificada.

\vspace{0.8em}
\leadq{\textquestiondown{}Realmente lo está?}

\vspace{0.3em}
Esa es la pregunta que abre este informe.
No vamos a optimizar la cartera. Vamos a practicarle una \emph{autopsia}:
abrirla capa por capa ---capital, dependencia, riesgo, posiciones efectivas, factores comunes---
hasta ver cuántas fuentes de riesgo independientes tiene realmente.

\figbig{fig01_portfolio_allocation_es.pdf}{La cartera bajo estudio: 60\% renta variable, 40\% renta fija.}

% ============================================================
% LAYER 1 — CAPITAL
% ============================================================
\newpage
\kicker{Capa 1 / Capital}
\section*{Primera mirada: el capital}

\leadq{\textquestiondown{}Cuántos activos tiene realmente esta cartera?}

Diez nombres en el extracto. Pero si el capital estuviera muy concentrado en pocos instrumentos,
ese número mentiría.

Medimos cuántas posiciones \emph{equiponderadas} producirían la misma dispersión de capital.
El resultado: \textbf{9,55}.

Casi diez. Hasta aquí, todo parece estar bien.

\insight{Desde el capital, la cartera está bien distribuida. El problema ---si existe--- no está en los pesos.}

\figbig{fig02_effective_positions_capital_es.pdf}{Diez activos nominales. 9,55 posiciones efectivas por capital.}

% ============================================================
% LAYER 2 — CORRELATION
% ============================================================
\newpage
\kicker{Capa 2 / Correlación}
\section*{Segunda mirada: cómo se mueven}

\leadq{\textquestiondown{}Los activos se comportan de manera independiente?}

Diez instrumentos no significan diez comportamientos distintos.
Si varios instrumentos reaccionan a las mismas noticias, están comprando ---en la práctica---
exposición repetida.

El mapa de correlaciones lo muestra de inmediato:
bloques rojos densos dentro de renta variable y dentro de renta fija;
casi independencia entre ambas clases.

\insight{Muchos instrumentos se mueven prácticamente al mismo tiempo. Contar nombres no es contar comportamientos.}

\figfull{fig03_correlation_heatmap_es.pdf}{Matriz de correlaciones de los retornos.}

\newpage
\kicker{Capa 2 / Correlación}
\section*{Los agrupamientos aparecen solos}

Cuando solo conectamos correlaciones fuertes, la red habla por sí sola:
un grupo de renta variable, un grupo de bonos, y TSLA como nodo aislado.

\insight{La cartera no tiene diez historias. Tiene unos pocos grupos de movimiento.}

\figfull{fig04_correlation_network_es.pdf}{Red de correlaciones fuertes ($\lvert\rho\rvert \ge 0{,}60$).}

% ============================================================
% LAYER 3 — RISK (Figura 5 protagonista)
% ============================================================
\newpage
\kicker{Capa 3 / Riesgo}
\section*{Tercera mirada: quién genera el riesgo}

\leadq{\textquestiondown{}Quién genera realmente el riesgo del portafolio?}

Hasta ahora miramos pesos y correlaciones.
Ahora la pregunta cambia: de cada peso invertido, \emph{cuánto riesgo sale}.

NVDA, TSLA y MU tienen el mismo peso que el resto del bloque de renta variable.
Su aporte al riesgo del portafolio es muy distinto.

\insight{Dos activos pueden tener exactamente el mismo peso y, aun así, aportar cantidades de riesgo completamente distintas.}

\newpage
\kicker{Capa 3 / Riesgo}
\section*{El cambio conceptual}

Aquí ocurre el giro: mismo capital, riesgo desigual.

\fighero{fig05_weight_vs_prc_es.pdf}{Peso en cartera vs.\ contribución al riesgo. La figura central del diagnóstico.}

\newpage
\kicker{Capa 3 / Riesgo}
\section*{60/40 de capital no es 60/40 de riesgo}

La etiqueta más famosa de la industria ---60/40--- describe cómo se reparte el \emph{dinero}.
No cómo se reparte el \emph{riesgo}.

\insight{El 40\% del capital explica apenas el 2,35\% del riesgo.}

\vspace{0.3em}
{\sffamily\small\color{Muted} La renta variable concentra el 97,65\% de la contribución total al riesgo.}

\figbig{fig06_capital_vs_risk_allocation_es.pdf}{Asignación de capital vs.\ asignación de riesgo.}

% ============================================================
% LAYER 4 — EFFECTIVE POSITIONS
% ============================================================
\newpage
\kicker{Capa 4 / Posiciones efectivas}
\section*{Cuarta mirada: posiciones que importan}

\leadq{\textquestiondown{}Cuántas posiciones efectivas existen realmente?}

Si aplicamos la misma lógica del capital ---cuántas posiciones equivalentes hay---
pero ahora sobre la distribución del riesgo, el número cae.

De 10 nominales a 9,55 por capital\ldots{} y a \textbf{6,58} por riesgo.

\insight{La cartera parece tener diez posiciones. Por riesgo, se comporta como si tuviera poco más de seis.}

\figbig{fig07_effective_positions_es.pdf}{De contar activos a medir riesgo: 10 $\rightarrow$ 9,55 $\rightarrow$ 6,58.}

% ============================================================
% LAYER 5 — PCA INTRO (página previa)
% ============================================================
\newpage
\kicker{Capa 5 / Factores comunes}
\section*{Antes de los componentes principales: una sola pregunta}

Llegamos a la última capa.
Antes de mirar gráficos, conviene fijar \emph{qué} estamos preguntando.

\vspace{1.2em}
\noindent{\sffamily\large\bfseries\color{Navy} \textquestiondown{}Qué intenta responder el análisis de componentes principales?}

\vspace{0.7em}
{\large\color{Ink}
No intenta reducir el número de activos.\\[0.35em]
No intenta optimizar la cartera.\\[0.35em]
No prescribe una asignación alternativa.}

\vspace{1.1em}
\noindent{\sffamily\large\bfseries\color{Navy} Intenta responder una única pregunta:}

\vspace{0.7em}
\begin{center}
{\LARGE\bfseries\color{Ink}
\textquestiondown{}Cuántas fuentes independientes\\[0.25em]
explican realmente el comportamiento\\[0.25em]
conjunto del portafolio?}
\end{center}

\vspace{1.4em}
\insight{Si pocas fuentes explican la mayor parte del movimiento conjunto, el número de instrumentos sobrestima la diversificación.}

\vspace{0.8em}
{\sffamily\small\color{Muted}%
En las páginas siguientes: gráfico de sedimentación, varianza acumulada y cargas factoriales.
Solo diagnóstico sobre la matriz de correlación.}

% ============================================================
% LAYER 5 — PCA RESULTS
% ============================================================
\newpage
\kicker{Capa 5 / Factores comunes}
\section*{Cuántas fuentes bastan}

Tres componentes explican al menos el 80\% de la variación conjunta.
Cinco alcanzan el 90\%.

Diez instrumentos. Pocas dimensiones dominantes.

\insight{El portafolio no necesita diez historias distintas para explicarse. Con tres ya se entiende la mayor parte.}

\figbig{fig08_pca_scree_es.pdf}{Gráfico de sedimentación: cuánta varianza aporta cada componente.}

\newpage
\kicker{Capa 5 / Factores comunes}
\section*{Qué representan esas fuentes}

La varianza acumulada confirma el umbral.
Las cargas factoriales muestran el patrón empírico ---sin imponer nombres a priori---:

\vspace{0.4em}
\begin{itemize}
  \item \textbf{PC1} ($\approx$49\%): se alinea con el bloque de renta variable.
  \item \textbf{PC2} ($\approx$27\%): se alinea con el bloque de renta fija.
  \item El resto: fracciones menores, sin etiqueta económica unívoca.
\end{itemize}

\figbig{fig09_pca_cumulative_es.pdf}{Varianza explicada acumulada.}

\newpage
\kicker{Capa 5 / Factores comunes}
\section*{Cargas factoriales: el mapa de cada fuente}

\insight{%
\textbf{PC1} representa principalmente renta variable.\\[0.2em]
\textbf{PC2} representa principalmente renta fija.\\[0.2em]
Los componentes restantes capturan variaciones más específicas.}

\figfull{fig10_pca_loadings_es.pdf}{Cargas factoriales escaladas (PC1--PC5).}

% ============================================================
% FINAL DIAGNOSIS — una sola página
% ============================================================
\newpage
\thispagestyle{fancy}
\kicker{Diagnóstico final}
\section*{Lo que quedó después de abrir la cartera}

\vspace{0.35em}
\begin{center}
{\sffamily\small\color{Muted} La cartera parecía tener}\\[0.35em]
{\fontsize{28}{32}\selectfont\bfseries\color{Ink} 10 activos}
\end{center}

\vspace{0.2em}
\begin{center}{\large\color{Muted}$\downarrow$}\end{center}
\vspace{0.15em}

\begin{center}
{\fontsize{22}{26}\selectfont\bfseries\color{Navy} 9,55}\\[0.05em]
{\sffamily\small\color{Muted} posiciones efectivas de capital}
\end{center}

\vspace{0.15em}
\begin{center}{\large\color{Muted}$\downarrow$}\end{center}
\vspace{0.15em}

\begin{center}
{\fontsize{22}{26}\selectfont\bfseries\color{Accent} 6,58}\\[0.05em]
{\sffamily\small\color{Muted} posiciones efectivas de riesgo}
\end{center}

\vspace{0.15em}
\begin{center}{\large\color{Muted}$\downarrow$}\end{center}
\vspace{0.15em}

\begin{center}
{\fontsize{28}{32}\selectfont\bfseries\color{Ink} 3}\\[0.05em]
{\sffamily\small\color{Muted} fuentes principales de riesgo}
\end{center}

\vspace{0.7em}
\noindent{\color{Navy}\rule{\textwidth}{2.5pt}}
\vspace{0.55em}

\begin{center}
{\LARGE\bfseries\color{Ink}
Un portafolio puede tener muchos activos\\[0.3em]
y muy pocas ideas.}
\end{center}

\vspace{0.45em}
\begin{center}
{\large\color{Slate}
Diez instrumentos no implican diez riesgos independientes.}
\end{center}

\vspace{0.55em}
\noindent{\color{Navy}\rule{\textwidth}{2.5pt}}

\vspace{0.5em}
{\sffamily\small\color{Muted}
Villar Capital --- Informe de investigación nro.\ 001 $|$ Diagnóstico, no optimización}

\end{document}
